% ============================================================
% Raccomandazioni finali
% ============================================================

\chapter{Raccomandazioni}
\label{cap:recommendations}

\section{Risposta alle domande guida}

\subsection{Meglio SOLIDWORKS o NX per droni?}

Per il nostro campo di applicazioni, **NX e' migliore come benchmark proprietario**. SOLIDWORKS resta molto competitivo per prototipazione meccanica e team piccoli, ma NX e' piu' adatto quando entrano superfici aerodinamiche complesse, compositi, model-based workflow, Simcenter e ottimizzazione automatica.

\subsection{STAR-CCM+ fa bene le mesh?}

Si. STAR-CCM+ ha una pipeline di meshing industriale robusta e integrata, con surface repair/wrapping, mesh polyhedral/trimmer, prism layers, controlli locali e adattamento~\cite{starccm_official,starccm_training}. Se si usa STAR-CCM+ come solver, il meshing interno e' la scelta consigliata. GMSH resta preferibile nella baseline open-source verso SU2/OpenFOAM.

\subsection{Come funziona HEEDS?}

HEEDS automatizza il processo CAD/CAE, lancia varianti, estrae risultati e applica algoritmi di esplorazione/ottimizzazione. A livello industriale permette di modificare parametri geometrici, condizioni operative, vincoli e obiettivi, poi di analizzare trade-off, Pareto front e sensitivita'. Il suo valore e' massimo quando ogni run CFD e' costosa e il team vuole esplorare molte piu' varianti di quante potrebbe testare manualmente.

\section{Ranking finale}

\begin{table}[H]
\centering
\caption{Ranking dei due casi studio proprietari}
\begin{tabular}{c P{5cm} c P{5.5cm}}
\toprule
\textbf{Rank} & \textbf{Workflow} & \textbf{Score} & \textbf{Uso consigliato} \\
\midrule
1 & NX + STAR-CCM+ + HEEDS & 4.20/5 & Benchmark industriale principale per droni/eVTOL complessi \\
2 & SOLIDWORKS + STAR-CCM+ + HEEDS & 3.90/5 & Alternativa pratica se il team e' gia' SOLIDWORKS-first \\
\bottomrule
\end{tabular}
\end{table}

\section{Strategia per la repo}

La repo dovrebbe mantenere due famiglie di workflow:

\begin{itemize}
    \item \textbf{Baseline replicabile}: OpenVSP + GMSH + SU2/OpenFOAM + Optuna, utile per didattica, automazione open e versionamento.
    \item \textbf{Benchmark industriale}: NX + STAR-CCM+ + HEEDS, utile per capire come lavorerebbe un team con licenze, HPC e processo proprietario maturo.
\end{itemize}

SOLIDWORKS + STAR-CCM+ + HEEDS va mantenuto nel catalogo come alternativa concreta, ma non come target ideale per il caso Airspeeder-like.

\section{Prossimi passi}

\begin{enumerate}
    \item Creare un micro-modello CAD comune: fusoliera semplificata, due bracci, fairing rotore e patch naming standard.
    \item Definire una matrice di validazione mesh in STAR-CCM+: coarse, medium, fine, y+ target, $C_D$ convergence.
    \item Creare una tabella HEEDS minima con 5 variabili, 2 vincoli e 2 obiettivi.
    \item Replicare lo stesso micro-caso nella baseline open source per confrontare costo, attrito operativo e riproducibilita'.
\end{enumerate}
